\section{Representations in categories}

A representation of a group is a way to view the group as acting on some mathematical
object while preserving the structure of that object.

For example, a group may act on:

\begin{itemize}
    \item a set;
    \item the set of roots of a polynomial;
    \item a vector space;
    \item a topological space;
    \item an object of a category.
\end{itemize}

The key idea is that a group action should preserve the underlying structure of the object.


Let \(\mathcal{C}\) be a category. Examples include:

\begin{itemize}
    \item \(\mathbf{Set}\), the category of sets;
    \item \(\mathbf{FinSet}\), the category of finite sets;
    \item \(\mathbf{Top}\), the category of topological spaces;
    \item \(\mathbf{Vect}_k\), the category of vector spaces over a field \(k\);
    \item \(\mathbf{vect}_k\), the category of finite-dimensional vector spaces over \(k\).
\end{itemize}

If \(X\) is an object of \(\mathcal{C}\), then the automorphism group of \(X\) is
\[
    \operatorname{Aut}_{\mathcal{C}}(X)
    =
    \{\, f:X\to X \mid f \text{ is an isomorphism in } \mathcal{C}\,\}.
\]
Equivalently, \(\operatorname{Aut}_{\mathcal{C}}(X)\) consists of all morphisms
\[
    f:X\to X
\]
which admit an inverse morphism.

\begin{definition}[Representation in a category]
    Let \(G\) be a group and let \(X\) be an object of a category \(\mathcal{C}\).
    A representation of \(G\) on \(X\) is a group homomorphism
    \[
        \rho:G\longrightarrow \operatorname{Aut}_{\mathcal{C}}(X).
    \]
\end{definition}

\begin{example}[Permutation representations]
    If \(\mathcal{C}=\mathbf{Set}\), then a representation of \(G\) on a set \(X\) is a
    group homomorphism
    \[
        \rho:G\longrightarrow \operatorname{Aut}_{\mathbf{Set}}(X).
    \]
    Since
    \[
        \operatorname{Aut}_{\mathbf{Set}}(X)=\operatorname{Sym}(X),
    \]
    this is precisely a permutation representation
    \[
        \rho:G\longrightarrow \operatorname{Sym}(X).
    \]
\end{example}

\begin{example}[Linear representations]
    If \(\mathcal{C}=\mathbf{Vect}_k\), then for a vector space \(V\) over \(k\),
    \[
        \operatorname{Aut}_{\mathbf{Vect}_k}(V)=\operatorname{GL}(V).
    \]
    Therefore a linear representation of \(G\) on \(V\) is a group homomorphism
    \[
        \rho:G\longrightarrow \operatorname{GL}(V).
    \]
\end{example}

\begin{example}[Topological representations]
    If \(\mathcal{C}=\mathbf{Top}\), then for a topological space \(X\),
    \[
        \operatorname{Aut}_{\mathbf{Top}}(X)=\operatorname{Homeo}(X),
    \]
    the group of homeomorphisms of \(X\). Thus a topological action of \(G\) on \(X\)
    is a homomorphism
    \[
        \rho:G\longrightarrow \operatorname{Homeo}(X).
    \]
\end{example}

\begin{definition}[\(\mathcal{C}\)-representation of a group]
    Let \(G\) be a group, and let \(\mathcal{C}\) be a category.
    A representation of \(G\) on \(\mathcal{C}\), or a
    \(\mathcal{C}\)-representation of \(G\), is a pair
    \[
        (X,\rho),
    \]
    where
    \[
        X\in \operatorname{Ob}(\mathcal{C})
    \]
    is an object of \(\mathcal{C}\), and
    \[
        \rho:G\longrightarrow \operatorname{Aut}_{\mathcal{C}}(X)
    \]
    is a group homomorphism.
\end{definition}

\begin{remark}
    If \((X,\rho)\) is a representation of \(G\), we also say that
    \(G\) acts on \(X\). For each \(g\in G\), the morphism
    \[
        \rho(g):X\longrightarrow X
    \]
    is an automorphism of \(X\) in the category \(\mathcal{C}\).
\end{remark}

\begin{definition}[Faithful representation]
    A representation \((X,\rho)\) of \(G\) is called faithful if
    \[
        \rho:G\longrightarrow \operatorname{Aut}_{\mathcal{C}}(X)
    \]
    is injective.
\end{definition}

\begin{definition}[Morphism of representations]
    Let \((X,\rho)\) and \((X',\rho')\) be two
    \(\mathcal{C}\)-representations of \(G\).
    A morphism
    \[
        f:(X,\rho)\longrightarrow (X',\rho')
    \]
    is a morphism
    \[
        f:X\longrightarrow X'
    \]
    in \(\mathcal{C}\) such that for every \(g\in G\), the following diagram
    commutes:
    \[
    \begin{tikzcd}
        X \arrow[r,"f"] \arrow[d,"\rho(g)"'] &
        X' \arrow[d,"\rho'(g)"] \\
        X \arrow[r,"f"'] &
        X'
    \end{tikzcd}
    \]
    Equivalently,
    \[
        f\circ \rho(g)=\rho'(g)\circ f
        \qquad
        \text{for all } g\in G.
    \]
\end{definition}

\begin{definition}[The category of representations]
    Let \(G\) be a group and let \(\mathcal{C}\) be a category.
    We define
    \[
        \operatorname{Rep}(G,\mathcal{C})
    \]
    to be the category whose objects are \(\mathcal{C}\)-representations
    of \(G\), and whose morphisms are morphisms of representations.

    More explicitly:
    \begin{itemize}
        \item an object of \(\operatorname{Rep}(G,\mathcal{C})\) is a pair
        \[
            (X,\rho),
        \]
        where \(X\in \operatorname{Ob}(\mathcal{C})\) and
        \[
            \rho:G\to \operatorname{Aut}_{\mathcal{C}}(X)
        \]
        is a group homomorphism;

        \item a morphism
        \[
            (X,\rho)\longrightarrow (X',\rho')
        \]
        is a morphism \(f:X\to X'\) in \(\mathcal{C}\) satisfying
        \[
            f\circ \rho(g)=\rho'(g)\circ f
            \qquad
            \text{for all } g\in G.
        \]
    \end{itemize}
\end{definition}

\begin{remark}
    The identity morphism of a representation \((X,\rho)\) is
    \[
        \operatorname{id}_X:X\to X,
    \]
    and the composition of morphisms is inherited from the category
    \(\mathcal{C}\). Hence \(\operatorname{Rep}(G,\mathcal{C})\) is itself
    a category.
\end{remark}
